\section{Background and Threat Model}
\label{sec:background}

\subsection{Honeypots and the Fidelity--Safety Tension}
\label{sec:tension}
A honeypot is a resource with no production value whose only purpose is to
be probed, attacked, and compromised~\cite{spitzner2002}. The literature
classifies honeypots by interaction depth~\cite{lopez2024}. Low-interaction
systems emulate a handful of protocols well enough to be scanned;
medium-interaction systems accept sessions and expose a bounded command
surface; high-interaction systems hand the attacker a real machine, with
the containment cost that implies. Production deployments overwhelmingly
favor low- and medium-interaction designs because they contain risk:
a scripted trap cannot be pivoted through, and its worst-case disclosure
is the script itself. Containment, however, is bought by scripting, and
the deception is a finite automaton that cannot answer a question it was
not compiled to answer; an attacker who strays off the script learns
nothing new and may learn the host is fake.

Deception research gives the tension a formal shape: a deception succeeds
when the attacker's belief about the system diverges, in a way the
operator benefits from, from the system's
reality~\cite{pawlick2021}. That framing exposes the catch-22
that confines classical designs. A honeypot is valuable exactly insofar
as the attacker believes it and engages with it, and engagement rewards
richness; richness, delivered by scripts, is where the risk lives,
because every scripted answer is a hand-written commitment about a
machine that does not exist. Classical practice manages this by keeping
the fiction poor. Stated in the adversary's terms, the trade is a
catch-22: the fiction must be rich to be believed and poor to be safe.

Raising interaction depth through generative models inverts the trade.
A generator can answer \emph{anything}, which is exactly the problem.
First-generation generative honeypots demonstrate that an LLM can sustain
a shell conversation, an LDAP directory, or an AI-agent lure convincingly
enough to hold modern attackers~\cite{valeros2024,reworr2024,jimenez2025},
and LLM-generated web deception~\cite{galah2024} extends the same idea to
HTTP surfaces. The field has since moved from feasibility to state
awareness: recent systems manage persistent state and validate generated
responses against state constraints~\cite{li2026llmsshh}, add guardrails
across web and API surfaces~\cite{laragutierrez2026genpot},
and price the design trilemma of flexibility, depth, and
deception~\cite{wang2026honeygpt}, with a recent systematic review
mapping the design space~\cite{bridges2026}. The question this paper
addresses is the complementary one those systems do not pose: what is
the \emph{security} status of a response in a deception system that
serves it? For a security tool, wrongness takes three forms, and all
three are silent failures. \emph{Inconsistency}: the fiction drifts, and
an attacker who cross-references two turns learns the host is synthetic.
\emph{Over-disclosure}: the model complies too well and emits credential
material that was never intended to be served. \emph{Unavailability}:
the generator is a rented service that can refuse, stall, or vanish
mid-session, and a deception endpoint that answers with a refusal has
told the attacker something true. Prior systems mitigate the first form
inside the generator's context window, but a validated response is
still not a \emph{certified} one: no verdict is recorded, no fallback
bounds what happens when validation fails, no gate decides whether
high-stakes bait may be served. The design in \S\ref{sec:design} treats
all three failure forms as first-class constraints with enforcing
mechanisms, and prices every mechanism those constraints admit against
cheaper certifiable alternatives -- decoupling the realism of the
fiction from the risk of the machinery that produces it by making every
generator output pass a check computed against a world the defender
fully controls.

\subsection{Threat Model}
\label{sec:threatmodel}
We consider an attacker who interacts with the honeypot over standard
remote-access protocols (SSH today, HTTP and Telnet supported by the same
machinery). The attacker may be unskilled (credential spraying, banner
grabbing) or expert (hands-on-keyboard, staged tooling, lateral-movement
reconnaissance); the behavioral difference between the two is what the
profiler must recover. The attacker's goals are those of the CTI
literature: initial access, credential harvest, privilege escalation, and
exfiltration, observable as techniques in the \attck matrix. Sessions are
attributed behaviorally, not cryptographically: the system may link
reconnaissance in one session to credential use in another, but an
attacker who rotates infrastructure presents as a new profile.

The attacker's \emph{observation model} is the served response and
nothing else. They cannot see the plan that produced it, the graph that
constrained it, or the verdicts that preceded it. Within that constraint
we assume arbitrary fingerprinting skill: the attacker may issue
conflicting commands, re-query facts stated many turns earlier, compare
notes across parallel sessions, and probe protocol edge cases
specifically to test whether the host is synthetic. Inconsistency is
therefore not a cosmetic defect but an information channel to the
adversary, which is why G2 (below) requires a zero-contradiction
condition rather than a low rate.

Timing is the one channel this observation model leaks. Response
latency is measurable: a session served in tens of seconds (the live
campaign's mean was ${\sim}44$~s, dominated by generative turns;
\S\ref{sec:discussion}), the fallback's fast vetted answers, and the
HITL gate's operator pause on exactly the highest-stakes turns are all
distinguishable by timing alone -- the serve-path provenance analysis
of \S\ref{sec:evaluation} separates LLM-served from deterministic rows
by latency alone. We treat latency homogenization as an open problem;
the deception properties this paper claims are those verifiable from
response content alone.

We do \emph{not} assume the attacker can compromise the honeypot's
control plane, escape the container sandbox, observe traffic between
honeypot components, or influence the retrieval corpus before a session
begins. Physical access, insider threat, and compromise of the defender's
workstation are out of scope.

Finally, we assume the attacker may be \emph{adaptive}: they may probe
systematically to map the fiction's boundaries, deliberately elicit
refusals to fingerprint the generator, re-attack across sessions to test
which facts persist, and seed inputs shaped to influence what the
retrieval engine surfaces. The design does not assume the deception
logic is secret; the fallback transcript library is assumed
disclosable, and the fiction's single-scope deployment means an attacker
who compromises one session's view learns one host, not the fleet. The
per-session scoping of the fiction graph (\S\ref{sec:graph}) means
cross-session facts -- including the observable difference between the
deception policies served to different attacker classes -- lie outside
the zero-contradiction guarantee, which is therefore strictly
within-session (Prop.~\ref{prop:consistency}); an attacker comparing
notes across sessions or probing the same fiction class across
re-attacks is testing exactly this boundary, and we treat cross-session
consistency as future work rather than a claim. What
the system does rely on keeping secret is small: the canary registry,
which must exist outside any view the attacker can reach, and the
serving-time verdicts of \S\ref{sec:design}. An adaptive adversary is
the reason the evaluation of \S\ref{sec:evaluation} measures the system
under targeted adversarial re-query rather than benign traffic.

We additionally consider the \emph{provider} as an involuntary third
party. The honeypot's generative components run against an inference API
whose safety alignment, rate limits, and availability are outside our
control. Across this project's deployments we observed the full failure
taxonomy firsthand: safety refusals of the deception task, budget
exhaustion under alignment pressure (86.7\% of generated turns on one
commercial provider), timeouts, transient network errors, and catalog
churn, in which a listed model disappears from the provider's offering
between one day and the next. A deployment must remain safe and useful
through every row of that taxonomy; \S\ref{sec:design} accordingly treats
the provider as an unreliable component, not a trusted oracle.

Defenders operate the system and are assumed to act within policy. The human-in-the-loop
gate exists precisely because some content classes (credential bait, for
instance) should not be served without a human decision, even when the
machinery would allow it.

\subsection{Vocabulary}
\label{sec:vocab}
Five terms carry precise meanings throughout the paper. A \emph{deception
plan} is the per-session specification the loop serves against: the facts
to project, the bait artifacts to deploy, and the attacker model they are
chosen for. The \emph{fiction graph} is the triple store binding every
projected host fact (hostname, kernel, service versions, credentials) to
one consistent world; \S\ref{sec:design} formalizes its consistency
property. A response is \emph{verified} only if a serving-time check
confirms every factual claim it makes against the graph; anything else
that must be sent is served from the \emph{deterministic fallback}, a
library of human-vetted transcripts. A \emph{honeytoken} is a deployed
artifact (a credential, a configuration value) whose value is a
canary-registered secret, so that any use of it elsewhere is attributable
to the disclosure event; honeytokens as deliberate, trackable
disclosures date to the earliest honeypot literature~\cite{spitzner2003}. \emph{TTP} denotes an ATT\&CK tactic or
technique~\cite{roy2023}; extraction of TTPs from attacker behavior is
the CTI product the system computes as a by-product of deceiving, and
retrieval-augmented grounding~\cite{lewis2020} over curated CTI corpora,
a technique adapted to security annotation in prior
work~\cite{lekssays2025}, is what supplies that product its evidence base.

\subsection{Design Goals}
\label{sec:goals}
The analysis above yields five goals. Each goal names the failure it
exists to prevent, the mechanism that enforces it (\S\ref{sec:design}),
and the measurement that falsifies it (\S\ref{sec:evaluation}):
\begin{itemize}
\item \textbf{G1 (Grounding).} Deception content should be conditioned on
who the attacker is and on current threat intelligence, not on a static
script. Prevents: irrelevant deception that no attacker class finds
plausible. Enforced by: the profiler and hybrid retrieval engine.
Measured by: the paired RAG-vs-no-RAG differential (RQ1).
\item \textbf{G2 (Consistency).} The fiction must never contradict itself.
Every artifact the attacker can observe (kernel version, listening ports,
databases, credentials) must resolve to one set of facts. Prevents: the
inconsistency channel of \S\ref{sec:tension}. Enforced by: the fiction
graph and the serving-time verifier. Measured by: contradiction rate
under adversarial re-query (RQ2).
\item \textbf{G3 (Safety).} No unverified content reaches the wire, and
dangerous content classes require explicit human approval. Prevents:
over-disclosure. Enforced by: the verifier, canary registration, and the
HITL gate. Measured by: verified-serve rate, danger-gate catches, and
HITL audit coverage (RQ2).
\item \textbf{G4 (Graceful degradation).} When generation, verification,
or the provider fails, the system must fall back to deterministic vetted
content rather than improvising. Prevents: unavailability of
\S\ref{sec:tension} from becoming disclosure or refusal.
Enforced by: the deterministic fallback path. Measured by: fallback
correctness and the A5 equivalence check (RQ2).
\item \textbf{G5 (Auditability).} Every accept, reject, and gate decision
must leave a persisted record detailed enough to reconstruct the session
after the fact. Prevents: unverifiable claims about the other four goals.
Enforced by: the audit trail. Measured by: audit-log coverage and the
replay determinism of \S\ref{sec:stability}.
\end{itemize}

One goal deserves a defense. G5 is not a logging feature; it makes the
other four falsifiable. A verification gate that does not
persist its rejections cannot be audited, only believed, and a deception
system that cannot be audited has no place in a security operation. The
evaluation of \S\ref{sec:evaluation} treats the audit trail as primary
evidence rather than instrumentation.
